\documentclass[12pt]{article}
\usepackage{amsmath,amssymb}
\usepackage{booktabs}
\usepackage{geometry}
\geometry{margin=1in}

\begin{document}

\section*{CSE 400 — Fundamentals of Probability in Computing}

\subsection*{Lecture 18 Scribe}

\textbf{Marginal PMF, Correlation, Covariance, Joint Gaussian Random Variables}

\section*{1. Marginal PMF}

\subsection*{1.1 Definition}

For discrete random variables $(X)$ and $(Y)$, the \textbf{joint PMF} is

\[
p_{XY}(x,y) = P(X=x, Y=y)
\]

The \textbf{marginal PMFs} are obtained by summing over the other variable:

\[
p_X(x) = \sum_y p_{XY}(x,y)
\]

\[
p_Y(y) = \sum_x p_{XY}(x,y)
\]

\textbf{Key Idea}

The marginal PMF of one variable is obtained by \textbf{summing the joint PMF over all possible values of the other variable}.

\subsection*{1.2 Marginal PMF from a Joint PMF Table}

Row sums give $p_X(x)$, and column sums give $p_Y(y)$.

\[
p_X(x_i) = \sum_j p_{ij}
\]

\[
p_Y(y_j) = \sum_i p_{ij}
\]

Normalization conditions:

\[
\sum_i p_X(x_i) = 1
\]

\[
\sum_j p_Y(y_j) = 1
\]

\subsection*{1.3 Example}

Let the joint PMF be

\[
\begin{array}{c|cc}
p_{XY}(x,y) & y=0 & y=1\\
\hline
x=0 & 1/8 & 1/4\\
x=1 & 3/8 & 1/4
\end{array}
\]

\textbf{Step-by-step computation}

\textbf{Compute $p_X(0)$}

\[
p_X(0) = \frac{1}{8} + \frac{1}{4}
\]

\[
= \frac{1}{8} + \frac{2}{8}
\]

\[
= \frac{3}{8}
\]

\textbf{Compute $p_X(1)$}

\[
p_X(1) = \frac{3}{8} + \frac{1}{4}
\]

\[
= \frac{3}{8} + \frac{2}{8}
\]

\[
= \frac{5}{8}
\]

\textbf{Compute $p_Y(0)$}

\[
p_Y(0) = \frac{1}{8} + \frac{3}{8}
\]

\[
= \frac{4}{8}
\]

\[
= \frac{1}{2}
\]

\textbf{Compute $p_Y(1)$}

\[
p_Y(1) = \frac{1}{4} + \frac{1}{4}
\]

\[
= \frac{1}{2}
\]

\textbf{Final Marginal PMFs}

\[
p_X(x)=
\begin{cases}
3/8 & x=0\\
5/8 & x=1
\end{cases}
\]

\[
p_Y(y)=
\begin{cases}
1/2 & y=0\\
1/2 & y=1
\end{cases}
\]

\section*{2. Correlation}

\subsection*{2.1 Definition}

A basic way to measure joint behavior is through the \textbf{product moment}.

In the continuous case,

\[
R_{XY} = E[XY]
\]

It tells us the \textbf{average value of the product $(XY)$}.

This is often called a \textbf{raw correlation measure}.

\subsection*{2.2 Properties}

\[
R_{XY} = E[XY]
\]

\begin{itemize}
\item Depends on the \textbf{scale of $(X)$ and $(Y)$}
\item Computed \textbf{about the origin}, not about the means
\end{itemize}

Therefore it mixes together

\begin{itemize}
\item mean values
\item spread
\item dependence
\end{itemize}

If

\[
R_{XY}=0
\]

then the variables are \textbf{orthogonal about the origin}.

To remove the effect of the means, we \textbf{center the variables first}, which leads to \textbf{covariance}.

\section*{3. Covariance}

\subsection*{3.1 Definition and Interpretation}

Covariance is defined as

\[
Cov(X,Y)=E[(X-\mu_X)(Y-\mu_Y)]
\]

where

\[
\mu_X = E[X], \quad \mu_Y = E[Y]
\]

\textbf{Interpretation}

First subtract the mean from each variable.

Then measure how the \textbf{centered quantities vary together}.

Covariance captures whether deviations from the means occur

\begin{itemize}
\item in the \textbf{same direction}
\item or in \textbf{opposite directions}
\end{itemize}

\textbf{Sign Interpretation}

\begin{itemize}
\item $Cov(X,Y) > 0$ $\rightarrow$ Positive covariance
\item $Cov(X,Y) < 0$ $\rightarrow$ Negative covariance
\item $Cov(X,Y) = 0$ $\rightarrow$ Zero covariance
\end{itemize}

Covariance tells the \textbf{direction of joint linear variation} between $(X)$ and $(Y)$.

However, the \textbf{magnitude depends on units} and cannot measure the strength of the relationship.

\subsection*{3.2 Algebraic Form}

Start with

\[
Cov(X,Y) = E[(X-\mu_X)(Y-\mu_Y)]
\]

Expand the product:

\[
=E[XY - X\mu_Y - \mu_XY + \mu_X\mu_Y]
\]

Apply linearity of expectation:

\[
=E[XY] - \mu_YE[X] - \mu_XE[Y] + \mu_X\mu_Y
\]

Since

\[
\mu_X = E[X], \quad \mu_Y = E[Y]
\]

Therefore

\[
Cov(X,Y)=E[XY]-\mu_X\mu_Y
\]

or

\[
Cov(X,Y)=E[XY]-E[X]E[Y]
\]

\subsection*{3.3 Properties of Covariance}

\textbf{1. Symmetry}

\[
Cov(X,Y)=Cov(Y,X)
\]

\textbf{2. Variance as a Special Case}

\[
Cov(X,X)=Var(X)
\]

\textbf{3. Linearity with Respect to Constants}

\[
Cov(aX+b,cY+d)=ac,Cov(X,Y)
\]

\textbf{4. Independence Implies Zero Covariance}

If

\[
X \perp Y
\]

then

\[
Cov(X,Y)=0
\]

However,

\textbf{zero covariance means uncorrelated, but it does not imply independence in general.}

\subsection*{3.4 Example 1}

Given

\[
Y=-2X+3
\]

Find

\[
Cov(X,Y)
\]

\textbf{Step 1: Use covariance formula}

\[
Cov(X,Y)=E[XY]-E[X]E[Y]
\]

\textbf{Step 2: Substitute expression for $(Y)$}

\[
Y=-2X+3
\]

\textbf{Step 3: Compute $(XY)$}

\[
XY=X(-2X+3)
\]

\[
=-2X^2+3X
\]

\textbf{Step 4: Compute expectation of $(Y)$}

\[
E[Y]=E[-2X+3]
\]

\[
=-2E[X]+3
\]

\textbf{Step 5: Substitute into covariance formula}

\[
Cov(X,Y)=E[-2X^2+3X]-E[X](-2E[X]+3)
\]

Expanding:

\[
=-2E[X^2]+3E[X]+2(E[X])^2-3E[X]
\]

\[
=-2E[X^2]+2(E[X])^2
\]

\textbf{Step 6: Use variance identity}

\[
Var(X)=E[X^2]-(E[X])^2
\]

Thus

\[
Cov(X,Y)=-2Var(X)
\]

\subsection*{3.5 Pearson’s Correlation Coefficient}

Pearson’s correlation coefficient measures the \textbf{effect of change in one variable when the other variable changes}.

\[
\rho_{XY}=\frac{Cov(X,Y)}{\sqrt{Var(X)Var(Y)}}
\]

or

\[
\rho_{XY}=\frac{Cov(X,Y)}{\sigma_X\sigma_Y}
\]

Properties:

\begin{itemize}
\item Measures \textbf{strength and direction of linear relationship}
\item \textbf{Normalized version of covariance}
\item \textbf{Unit-free}
\end{itemize}

Range:

\[
-1 \le \rho_{XY} \le 1
\]

Interpretation:

\begin{itemize}
\item $\rho_{XY}>0$ $\rightarrow$ positive linear trend
\item $\rho_{XY}<0$ $\rightarrow$ negative linear trend
\item $\rho_{XY}=0$ $\rightarrow$ no linear correlation
\end{itemize}

\section*{4. Jointly Gaussian Random Variables}

\subsection*{4.1 Definition}

A pair $((X,Y))$ is \textbf{jointly Gaussian} if

\begin{itemize}
\item its \textbf{joint density has the bivariate Gaussian form}
\item the \textbf{marginals are Gaussian}
\end{itemize}

\subsection*{4.2 Algebraic Form}

The \textbf{joint Gaussian PDF} is

\[
f_{XY}(x,y)=
\frac{1}{2\pi\sigma_X\sigma_Y\sqrt{1-\rho^2}}
\exp
\left(
-\frac{1}{2(1-\rho^2)}
\left[
\left(\frac{x-\mu_X}{\sigma_X}\right)^2
+
\left(\frac{y-\mu_Y}{\sigma_Y}\right)^2
-
2\rho
\left(\frac{x-\mu_X}{\sigma_X}\right)
\left(\frac{y-\mu_Y}{\sigma_Y}\right)
\right]
\right)
\]

Where

\begin{itemize}
\item $( \mu_X, \mu_Y )$ = means
\item $( \sigma_X, \sigma_Y )$ = standard deviations
\item $( \rho )$ = Pearson correlation coefficient
\end{itemize}

\subsection*{4.3 Interpretation}

Example from the lecture:

Relationship between

\begin{itemize}
\item $(X)$: Atmospheric CO$_2$ concentration (ppm)
\item $(Y)$: Global surface temperature anomaly ($^\circ$C)
\end{itemize}

Because both variables vary due to many random environmental processes such as

\begin{itemize}
\item solar radiation
\item ocean circulation
\item emissions
\item weather variability
\end{itemize}

their joint behavior can often be approximated by a \textbf{Joint Gaussian Distribution}.

\end{document}